\section{Discussion}

\subsection{Architectural Generality}

三层解耦设计的核心价值在于职责分离——路由层可独立替换分类策略，角色层可独立修改人设，工作引擎可独立扩展子图。这种分层可复用到其他 Agent 场景（如客服、教育、运维），只需替换各层的具体实现。例如，在客服场景中，路由层可替换为基于 BERT 的意图分类器，角色层可替换为客服话术模板，工作引擎可扩展工单处理子图。

\subsection{Engineering Trade-offs}

\textbf{关键词门控 vs LLM 意图分类}：前者零延迟、零成本，但召回有限（类似 Rasa~\cite{bocklisch2017rasa} 的规则方法）；后者精度更高但增加延迟和 token 消耗（如 GPT-3~\cite{brown2020gpt3} 的零样本分类）。当前系统选择关键词门控作为默认路径，将 LLM 调用推迟到参数抽取阶段，是延迟-精度-成本的三角权衡。

\textbf{视觉感知在无 GPU 桌面端的可行性}：YOLOv8~\cite{jocher2023yolov8} ONNX CPU 推理（12 MB 模型）在普通笔记本上可在 1--2 秒内完成单次推理，60 秒的轮询间隔确保不干扰用户操作。

\subsection{Current Limitations}

本系统存在以下局限：（1）记忆系统无持久化，压缩归档和工作记忆层未实现；（2）视觉检测仅 6 类 UI 元素，无法理解复杂 UI 语义；（3）代码执行无沙箱隔离（本地子进程）；（4）情绪系统为全局单例，非 per-session。

\subsection{Positioning vs. Related Work}

Claude Computer Use~\cite{anthropic2024computeruse} 等系统具备更强的视觉理解和通用操作能力，但依赖云端大模型和 GPU 推理。本系统的优势在于轻量（全栈可在普通桌面端运行）、可离线（视觉推理仅需 CPU）、角色个性化（人设驱动的交互体验），适合低资源桌面场景下的陪伴型 AI 助手定位。

\subsection{Future Work}

未来工作方向包括：（1）补全三层记忆架构中的压缩归档与磁盘持久化；（2）扩展 YOLOv8 类别至代码编辑器、终端等桌面应用特定元素；（3）引入沙箱执行环境（如 Docker 容器）提升代码执行安全性；（4）将情绪系统改为 per-session 以支持多用户场景；（5）增加 LLM 意图分类作为可选的高精度路径，实现关键词门控与 LLM 的动态路由。
